\chapter{Hypotension}

\section*{Definition}
Hypotension is defined as systolic blood pressure less than 90 mmHg or mean arterial pressure less than 65 mmHg in adults, or a drop in systolic BP of more than 20 mmHg from baseline. It is clinically significant when associated with symptoms of hypoperfusion (dizziness, syncope, confusion, oliguria).

\section*{What Happens in Your Body}
Hypotension results from reduced cardiac output, decreased systemic vascular resistance, or a combination of both. The compensatory baroreceptor reflex increases heart rate, contractility, and peripheral narrowing of blood vessels; hypotension ensues when these mechanisms are overwhelmed or impaired.

\section*{Causes}
\begin{itemize}
  \item \textbf{Hypovolemic:} Hemorrhage, dehydration, burns, diarrhea, vomiting, excessive diuresis
  \item \textbf{Cardiogenic:} Acute myocardial infarction, myocarditis, arrhythmia, heart failure, valvular dysfunction, cardiac tamponade
  \item \textbf{Distributive:} Sepsis (most common cause in hospitalized people), anaphylaxis, neurogenic shock (spinal cord injury), adrenal insufficiency
  \item \textbf{Orthostatic:} Volume depletion, autonomic neuropathy (diabetes, Parkinson disease), prolonged bed rest, medications (alpha-blockers, antidepressants)
  \item \textbf{Other:} Medications (antihypertensives, vasodilators), pregnancy (compression of IVC), postprandial hypotension, adrenal crisis
\end{itemize}

\section*{When to See a Doctor}
See your doctor if:
\begin{itemize}
  \item Systolic BP less than 90 mmHg with symptoms of dizziness, confusion, or syncope
  \item Signs of shock: pale, cool, clammy skin, tachycardia, diminished pulses, altered mental status, oliguria
  \item Acute blood loss (hematemesis, melena, trauma) with hypotension
  \item Hypotension following a new medication dose or allergic reaction
  \item Orthostatic hypotension with recurrent syncope or falls
\end{itemize}

\section*{Related Symptoms}
\begin{itemize}
  \item Dizziness, lightheadedness, syncope, blurred vision, fatigue, confusion, nausea, cold extremities, oliguria
\end{itemize}
\textbf{Important:} Orthostatic hypotension is confirmed by a drop in systolic BP of at least 20 mmHg or diastolic BP of at least 10 mmHg within 3 minutes of standing from a supine position.

\section*{Self-Care / Home Management}
\begin{itemize}
  \item Increase fluid and sodium intake if no contraindications (heart failure, renal disease)
  \item Rise slowly from sitting or lying positions to avoid orthostatic drops
  \item Wear compression stockings to reduce venous pooling in the lower extremities
  \item Avoid alcohol, prolonged standing, and large carbohydrate-heavy meals (postprandial hypotension)
  \item Elevate the head of the bed to reduce nocturnal diuresis and improve morning blood pressure
\end{itemize}

\section*{How Your Doctor Will Evaluate You}
\begin{itemize}
  \item Orthostatic vital signs: supine, sitting, and standing (1 and 3 minutes) blood pressure and heart rate
  \item ECG for cardiac ischemia or arrhythmia
  \item Laboratory: CBC (anemia, hemorrhage), basic metabolic panel (dehydration, renal function), lactate (tissue hypoperfusion), troponin, blood cultures if sepsis suspected
  \item Bedside ultrasound (RUSH protocol): assess cardiac function, IVC collapsibility, pneumothorax, free fluid
  \item Autonomic testing for suspected neurogenic orthostatic hypotension (tilt-table test)
\end{itemize}

\section*{Treatment Options}
\begin{itemize}
  \item IV fluid resuscitation (1--2 L crystalloid) for hypovolemic and distributive shock
  \item Blood transfusion for hemorrhagic shock; address source of bleeding
  \item Vasopressors (norepinephrine first-line) for fluid-refractory septic or cardiogenic shock
  \item Inotropic support (dobutamine) for cardiogenic shock with low cardiac output
  \item Remove or adjust offending medications; treat underlying cause (antibiotics for sepsis, steroids for adrenal insufficiency)
\end{itemize}

\section*{References}
\begin{thebibliography}{99}
\bibitem{ref1} Freeman R, Wieling W, Axelrod FB, et al. "Consensus Statement on the Definition of Orthostatic Hypotension, Neurally Mediated Syncope and the Postural Tachycardia Syndrome." Auton Neurosci. 2011;161(1--2):46--48.
\bibitem{ref2} Gupta V, Lipsitz LA. "Orthostatic Hypotension in the Elderly: Diagnosis and Treatment." Am J Med. 2007;120(10):841--847.
\bibitem{ref3} Low PA, Tomalia VA. "Orthostatic Hypotension: Mechanisms, Causes, Management." J Clin Neurol. 2015;11(3):220--226.
\bibitem{ref4} Lahrmann H, Cortelli P, Hilz M, et al. "EFNS Guidelines on the Diagnosis and Management of Orthostatic Hypotension." Eur J Neurol. 2006;13(9):930--936.
\bibitem{ref5} Shibao C, Grijalva CG, Raj SR, et al. "Orthostatic Hypotension and Associated Risk of Hospitalization and Mortality." J Am Geriatr Soc. 2007;55(10):1503--1509.
\bibitem{ref6} Kasper DL, Fauci AS, Hauser SL, et al. "Harrison's Principles of Internal Medicine." 21st ed. McGraw-Hill; 2022.
\end{thebibliography}
\clearpage
